\prefacesection{Abstract}
\justifying
Modern software teams increasingly turn to large language models (LLMs) to accelerate development, yet many tools still generate code in an ad-hoc, one-shot fashion with weak structure, limited verification, and poor support for iterative refinement. This project, \textbf{MetaGPT}, implements an end-to-end agentic framework that transforms natural language prompts into complete software projects using a coordinated pipeline of specialized LLM-driven agents. The system enables users to create projects, inspect intermediate artifacts, and iteratively refine outputs through an interactive interface.

MetaGPT models a software team through four core agents---Manager, Architect, Engineer, and QA---each governed by explicit Standard Operating Procedures (SOPs). The Manager converts informal prompts into structured requirements, the Architect designs system architecture and file layouts, the Engineer generates implementation code, and the QA agent derives tests and quality feedback. These agents are executed as a stateful workflow with preserved intermediate results, enabling traceability from requirements to design to implementation. The system persists project state and generated artifacts to support repeatability, recovery, and long-running iterations, and it incorporates retrieval utilities to ground future improvements in existing project context. Evaluation on representative software prompts shows that this SOP-driven, workflow-orchestrated approach yields more consistent project structures, improved observability of agent outputs, and better support for iterative refinement compared to naive, single-step code generation, demonstrating MetaGPT’s viability as a practical foundation for agentic software engineering.
\\[0.5cm]